\section{Evaluation}
\label{sec:evaluation}

We evaluate three engines --- the KG baseline, the KG-mem storage-layer
control, and SFDB --- on the same fact streams, the same query set, and
the same verification harness. Every benchmark result reported below has
passed three-way cross-engine semantic equivalence checking
(Section~\ref{sec:correctness}); rows for which any two engines produced
divergent canonical result sets would be flagged and excluded, and no such
divergence occurred in the runs reported here
(Section~\ref{sec:eval:verification}).

\subsection{Experimental Setup}
\label{sec:eval:setup}

The measurements were taken on a single Windows~11 machine with an
\texttt{\showcpu}\ processor, \texttt{\showmemory}\ of memory, and Python
$3.14$. Each latency figure is the arithmetic mean of ten timed runs
following two warm-up iterations, all on the same fact stream and
query set, and every latency table reports the $95\%$ confidence
half-width alongside each mean. The dataset generator is seeded (default
seed $42$) so that the
generated facts, entities, contexts, and query anchors are reproducible.
Configuration, git commit, and \texttt{uv.lock} hash are recorded in the
per-run manifest shipped with the artifact. Every number in this section
--- including the headline ranges quoted in the abstract and
introduction --- is generated directly from a benchmark run by
\texttt{scripts/generate\_tables.py}; none is hand-transcribed
(Section~\ref{sec:artifact:reproducibility}).

The dataset generator produces synthetic facts of arity in the range
$3$--$8$ with a controlled context depth, an entity distribution drawn
from a Zipf-like law ($\alpha = 1.2$) that mirrors the degree skew typical
of real knowledge graph data, and a temporal validity envelope on $40\%$ of
facts. We report results at \benchscalecount{} scales: \benchscales{}
facts. At the largest scale the KG representation holds roughly $1.5$
million reified triples. We do not report results at $10^6$ facts, and
the paper does not extrapolate to that scale.

The three engines are configured as follows. The \emph{KG baseline}
persists dictionary-encoded reified triples to in-memory SQLite and
answers queries through SQLite (Section~\ref{sec:arch:kg}). The
\emph{KG-mem control} runs the identical reification and query logic with
the SQLite dictionary tables and triple indexes replaced by Python dicts;
it exists to separate the effect of context-indexed storage from the
effect of answering queries without SQLite round trips
(Section~\ref{sec:eval:control}). \emph{SFDB} stores each fact atomically
and answers queries from its in-process stalk, neighbourhood, context,
and temporal indexes, writing through to SQLite on insert only. One
consequence is stated up front: KG-mem performs no write-through
persistence, so its \emph{insert} latencies are not comparable to the two
write-through engines and are reported for completeness only; KG-mem is a
query-path control.

We classify queries into five families whose performance profiles are
fundamentally different:

\begin{description}
    \item[LOOKUP:] retrieve every fact anchored at a specific entity (subject
          match). These are the queries for which stalk indexing is designed.
    \item[GLOBAL:] full-scan queries with no context or entity anchor.
          These force the engine to touch every stored fact.
    \item[CONTEXT:] retrieve every fact stated at a given context or any of
          its descendants in the context poset --- the query shape the
          context-indexed design is named for, and the one the paper's own
          motivating example describes (``every fact that holds in the
          \texttt{clinical.adult} context''). The anchor context used here,
          \texttt{ctx.0}, covers roughly $45\%$ of the generated facts at
          every scale, a genuine sub-tree rather than a single leaf or the
          whole graph.
    \item[TEMPORAL (bounded):] range queries over the temporal envelope of a
          fact, bounded on both ends (a one-year window in the dataset's
          temporal span). SFDB narrows these via its year-bucketed
          temporal index; the KG variants have no temporal index and
          always scan and filter reconstructed facts
          (Section~\ref{sec:complexity}).
    \item[TEMPORAL (unbounded):] range queries with a start bound and no end
          bound (``everything from $2023$-$01$-$01$ onward''). SFDB resolves
          these against a second, flat, binary-searchable temporal index
          rather than the year-bucketed one used for the bounded case;
          Section~\ref{sec:discussion:temporal-unbounded} describes why the
          bucketed index alone cannot answer this case. An earlier
          revision of the suite passed a bare year here and was silently
          re-bounded to a single-year window by a shared parsing
          shorthand; Section~\ref{sec:discussion:threats} records that
          bug, and a regression test now pins this query to a genuinely
          open end.
\end{description}

\subsection{Insert Throughput}
\label{sec:eval:insert}

SFDB's insert throughput relative to the KG baseline does not hold a
single direction across scales (Table~\ref{tab:insert}): at $100$ facts
KG is faster ($8.24$\,ms vs.\ SFDB's $10.09$\,ms, a $1.22\times$ margin),
while at $1\,000$, $10\,000$, and $100\,000$ facts SFDB is faster by
$1.3$--$1.5\times$. We report the crossover rather than average over it:
SFDB stores the fact atomically as
a single section, while the KG engine must decompose the fact into $k + 4$
triples and update three indexes per triple; both write through to
SQLite. At the smallest scale the fixed per-insert cost of registering a
fact into eight to ten open sets (Section~\ref{sec:eval:memory}) is not
yet amortised against that per-triple decomposition cost, and the two
engines' $95\%$ confidence intervals overlap narrowly rather than
separating cleanly, so we do not read the $100$-fact point as a strong
claim in either direction --- only as evidence that SFDB's insert
advantage is not unconditional. KG-mem, which skips persistence entirely, is faster than both ---
included in Table~\ref{tab:insert} for completeness, not comparison.

% Auto-generated by scripts/generate_tables.py from results/paper_suite_summary.json
% Do not edit by hand.
\begin{table}[htbp]
\centering
\caption{Insert latency in ms (mean over ten runs $\pm$ 95\% CI
half-width). Lower is better. KG-mem is the storage-layer control; it
performs no write-through persistence, so its insert cost is not
comparable to the two write-through engines (see text). The ratio
column is SFDB\,/\,KG.}
\label{tab:insert}
\begin{tabular}{lrrrr}
\toprule
Facts & KG (ms) & KG-mem (ms) & SFDB (ms) & SFDB/KG \\
\midrule
$100$ & $11.55 \pm 0.47$ & $3.49 \pm 0.23$ & $12.02 \pm 0.31$ & $1.04\times$ \\
$1,000$ & $133.8 \pm 1.7$ & $37.59 \pm 1.0$ & $104.7 \pm 1.9$ & $0.78\times$ \\
$10,000$ & $1,535.5 \pm 21.4$ & $370.9 \pm 9.7$ & $1,256.6 \pm 25.0$ & $0.82\times$ \\
$100,000$ & $14,474.8 \pm 1,009$ & $3,122.4 \pm 494$ & $13,601.9 \pm 2,018$ & $0.94\times$ \\
\bottomrule
\end{tabular}
\end{table}


\begin{figure}[htbp]
\centering
\includegraphics[width=0.6\textwidth]{figures/throughput.pdf}
\caption{Insert latency (SFDB/KG ratio annotated). KG-mem performs no
write-through persistence and is shown for completeness only.}
\label{fig:throughput}
\end{figure}

\subsection{Query Performance by Class}
\label{sec:eval:query}

Table~\ref{tab:query} reports LOOKUP, GLOBAL, CONTEXT, bounded TEMPORAL, and
unbounded TEMPORAL means with $95\%$ confidence half-widths at every
scale, for all three engines, with SFDB's ratio against both KG variants.

% Auto-generated by scripts/generate_tables.py from results/paper_suite_summary.json
% Do not edit by hand.
\begin{table}[htbp]
\centering
\caption{Query latency in ms by class and scale (mean over ten runs
$\pm$ 95\% CI half-width). Lower is better. KG-mem is the
storage-layer control: identical reification and query logic to KG with
dict indexes in place of SQLite. Ratios are SFDB\,/\,KG and
SFDB\,/\,KG-mem; below $1$ means SFDB is faster.}
\label{tab:query}
\small
\begin{tabular}{llrrrrr}
\toprule
Class & Facts & KG (ms) & KG-mem (ms) & SFDB (ms) & vs KG & vs KG-mem \\
\midrule
LOOKUP & $100$ & $0.9251 \pm 0.04$ & $0.3316 \pm 0.01$ & $0.0411 \pm 0.004$ & $0.04\times$ & $0.12\times$ \\
LOOKUP & $1,000$ & $9.76 \pm 0.07$ & $4.82 \pm 0.09$ & $0.3062 \pm 0.02$ & $0.03\times$ & $0.06\times$ \\
LOOKUP & $10,000$ & $95.03 \pm 1.2$ & $52.67 \pm 0.95$ & $3.34 \pm 0.10$ & $0.04\times$ & $0.06\times$ \\
LOOKUP & $100,000$ & $1,125.6 \pm 17.6$ & $744.8 \pm 11.2$ & $57.66 \pm 0.71$ & $0.05\times$ & $0.08\times$ \\
\midrule
GLOBAL & $100$ & $2.53 \pm 0.05$ & $0.8731 \pm 0.02$ & $0.0465 \pm 0.001$ & $0.02\times$ & $0.05\times$ \\
GLOBAL & $1,000$ & $28.27 \pm 0.92$ & $11.06 \pm 0.11$ & $0.2869 \pm 0.02$ & $0.01\times$ & $0.03\times$ \\
GLOBAL & $10,000$ & $302.9 \pm 8.0$ & $132.4 \pm 3.8$ & $4.03 \pm 0.62$ & $0.01\times$ & $0.03\times$ \\
GLOBAL & $100,000$ & $3,310.0 \pm 49.5$ & $1,599.0 \pm 33.1$ & $61.64 \pm 1.8$ & $0.02\times$ & $0.04\times$ \\
\midrule
CONTEXT & $100$ & $1.27 \pm 0.09$ & $0.4450 \pm 0.05$ & $0.0927 \pm 0.001$ & $0.07\times$ & $0.21\times$ \\
CONTEXT & $1,000$ & $14.44 \pm 0.07$ & $6.23 \pm 0.13$ & $0.6933 \pm 0.02$ & $0.05\times$ & $0.11\times$ \\
CONTEXT & $10,000$ & $165.3 \pm 1.8$ & $79.88 \pm 3.0$ & $34.53 \pm 0.46$ & $0.21\times$ & $0.43\times$ \\
CONTEXT & $100,000$ & $1,946.1 \pm 29.8$ & $1,041.6 \pm 18.7$ & $656.3 \pm 34.4$ & $0.34\times$ & $0.63\times$ \\
\midrule
TEMPORAL (bounded) & $100$ & $0.6068 \pm 0.02$ & $0.1990 \pm 0.003$ & $0.4952 \pm 0.005$ & $0.82\times$ & $2.49\times$ \\
TEMPORAL (bounded) & $1,000$ & $8.68 \pm 0.99$ & $3.80 \pm 0.05$ & $5.45 \pm 0.35$ & $0.63\times$ & $1.43\times$ \\
TEMPORAL (bounded) & $10,000$ & $96.52 \pm 10.0$ & $46.49 \pm 0.83$ & $58.11 \pm 0.83$ & $0.60\times$ & $1.25\times$ \\
TEMPORAL (bounded) & $100,000$ & $1,128.5 \pm 73.0$ & $724.2 \pm 11.4$ & $814.4 \pm 19.8$ & $0.72\times$ & $1.12\times$ \\
\midrule
TEMPORAL (unbounded) & $100$ & $1.02 \pm 0.04$ & $0.3350 \pm 0.007$ & $0.5724 \pm 0.05$ & $0.56\times$ & $1.71\times$ \\
TEMPORAL (unbounded) & $1,000$ & $11.41 \pm 0.36$ & $5.17 \pm 0.05$ & $5.26 \pm 0.07$ & $0.46\times$ & $1.02\times$ \\
TEMPORAL (unbounded) & $10,000$ & $127.8 \pm 8.2$ & $63.23 \pm 1.6$ & $59.91 \pm 1.3$ & $0.47\times$ & $0.95\times$ \\
TEMPORAL (unbounded) & $100,000$ & $1,509.8 \pm 20.0$ & $902.6 \pm 20.1$ & $835.2 \pm 9.5$ & $0.55\times$ & $0.93\times$ \\
\bottomrule
\end{tabular}
\end{table}


\begin{figure}[htbp]
\centering
\includegraphics[width=\textwidth]{figures/latency.pdf}
\caption{Query latency across scales, log axis, all three engines.}
\label{fig:latency}
\end{figure}

Against the KG baseline, SFDB is \rangelookupkg{} faster on LOOKUP
and \rangeglobalkg{} faster on GLOBAL. On
LOOKUP, SFDB is faster because the stalk index returns a
complete fact for the anchor entity in a single hash lookup, while the KG
engine must locate every triple that mentions the entity through its POS
index, gather them by event identifier, and reassemble the fact through
joins. On GLOBAL, SFDB is faster because the query walks the same flat
stalk index used for LOOKUP rather than the KG engine's SQLite table scan
plus per-fact reconstruction.

On CONTEXT, SFDB is \rangecontextkg{} faster than the KG baseline: every
fact is registered, at insert time, into the open set of its own context
and every ancestor context's open set (Section~\ref{sec:arch:sheaf}), so
the open set named for the query's anchor context already holds the full
sub-tree union before the query runs and the query is a single lookup
into it, while KG evaluates an exact-match-or-prefix
filter over its context-indexed triple table and still pays reified-triple
reconstruction per matching fact. This is the query shape the context
poset exists to answer, and it is measured here for the first time in
this evaluation's history --- an earlier revision of this paper reported
LOOKUP, GLOBAL, and TEMPORAL numbers without ever benchmarking the one
query class its own title names.

On the TEMPORAL classes the picture against the KG baseline is
narrower --- \rangetemporalkg{} on bounded ranges and
\rangetemporalunbkg{} on unbounded ranges --- because both engines do
comparable per-candidate interval filtering once the candidate set is
identified; SFDB's advantage here is in narrowing the candidate set
(year buckets for the bounded case, a binary search into the flat
start/end index for the unbounded case) rather than in reconstruction.
Section~\ref{sec:eval:control} shows how much of each advantage is
attributable to the storage layer rather than to context-indexed
storage itself.

\begin{figure}[htbp]
\centering
\includegraphics[width=0.6\textwidth]{figures/speedup.pdf}
\caption{SFDB/KG latency ratio by query class; below the parity
line means SFDB is faster.}
\label{fig:speedup}
\end{figure}

\subsection{The KG-mem Storage-Layer Control}
\label{sec:eval:control}

The comparison between SFDB and the KG baseline conflates two variables:
SFDB both stores facts atomically under a context poset \emph{and}
answers queries from in-process dict indexes, while the KG baseline both
reifies facts \emph{and} answers queries through SQLite. The KG-mem
control isolates the second variable: it reifies and reconstructs
exactly as the KG baseline does, but its dictionary tables and triple
indexes are Python dicts, so the SFDB-vs-KG-mem ratio measures what
context-indexed atomic storage buys once the storage-layer difference is
removed.

The result decomposes the headline speedups cleanly. On LOOKUP, SFDB
remains \rangelookupkgmem{} faster than KG-mem: most of the LOOKUP
advantage is attributable to avoiding reification --- returning a stored
fact instead of reassembling it from triples --- not to avoiding SQLite.
On GLOBAL the same holds at \rangeglobalkgmem{}. On CONTEXT, the ratio
against KG-mem is \rangecontextkgmem{}: both engines resolve the query in
time proportional to the result size, not to $N$, so this is not an
asymptotic gap but a constant-factor one --- SFDB's context-to-open-set
index resolves the anchor context to a single pre-populated open set
(every fact is registered into its own context's open set and every
ancestor's, at insert time, so the sub-tree union already exists before
the query runs), while KG-mem enumerates the handful of distinct context
buckets under the anchor and unions their fact-id sets. This is why
CONTEXT, unlike LOOKUP and GLOBAL, is a genuine test of whether that
constant factor is large enough to matter, rather than a foregone
conclusion.
On the TEMPORAL
classes, by contrast, the picture reverses at most scales rather than
merely narrowing. On bounded TEMPORAL, SFDB is slower than KG-mem at
three of the four scales measured ($100$, $10\,000$, and $100\,000$
facts) and marginally faster only at $1\,000$
(\rangetemporalkgmem{} across the full range, which straddles parity
rather than sitting above it); on unbounded TEMPORAL the pattern is
mixed, slower at the smallest scale and faster or near parity at the
three larger ones (\rangetemporalunbkgmem{}). The TEMPORAL
advantages over the SQLite-backed baseline reported earlier in this
section are therefore not evidence of a context-indexed design win at
all: they are storage-layer effects that a dict-indexed reified store
would recover just as well, and once that effect is controlled for,
context-indexed storage shows no consistent advantage on either
TEMPORAL class --- if anything a mild disadvantage on the bounded one.
We report this as a negative result rather than reframe it, consistent
with the paper's house style (Section~\ref{sec:discussion:negative}).
The honest summary is that context-indexed
storage pays for retrieval shaped like ``give me whole facts'' or ``give
me this context's sub-tree,'' and does not
meaningfully pay for range scans whose cost is dominated by
per-candidate filtering both designs share.

\begin{figure}[htbp]
\centering
\includegraphics[width=\textwidth]{figures/control.pdf}
\caption{Storage-adjusted speedup: SFDB/KG-mem latency ratio by query
class. Below the parity line means SFDB is faster than the dict-indexed
control.}
\label{fig:control}
\end{figure}

\subsection{External Reference Point: rdflib}
\label{sec:eval:rdflib}

Every comparison above is against engines we wrote ourselves. As an
external check, we run all five query classes against
rdflib~\cite{rdflib} --- a real, independent, pure-Python in-memory RDF
library --- at $100$, $1\,000$, and $10\,000$ facts, driven through
SPARQL text with the same overlap semantics the typed queries implement,
and with result counts checked against SFDB on every class at every
scale (Table~\ref{tab:rdflib}). The scope of this comparison is stated
plainly: rdflib is not a disk-backed production triple store, its
latencies include SPARQL parsing on every run (which is how the library
is normally driven), and we cap it at $10^4$ facts because its SPARQL
engine does not sustain the scan-shaped classes at $10^5$ in reasonable
benchmark time. It is an independent implementation by other authors,
which the in-house baselines are not, and every reported row returned
identical result counts to SFDB.

% Auto-generated by scripts/generate_tables.py from results/rdflib_reference.json
% Do not edit by hand.
\begin{table}[htbp]
\centering
\caption{External reference point: rdflib (a real, independent,
pure-Python in-memory RDF library queried via SPARQL) versus SFDB on
the same fact streams (mean over ten runs). The ratio is
SFDB\,/\,rdflib; below $1$ means SFDB is faster. The last column
records whether both systems returned the same number of results.
rdflib is not a disk-backed production triple store; see the text for
scope.}
\label{tab:rdflib}
\small
\begin{tabular}{llrrrc}
\toprule
Class & Facts & rdflib (ms) & SFDB (ms) & Ratio & Counts \\
\midrule
LOOKUP & $100$ & $1.17$ & $0.0236$ & $0.02\times$ & \checkmark \\
LOOKUP & $1,000$ & $3.25$ & $0.2241$ & $0.07\times$ & \checkmark \\
LOOKUP & $10,000$ & $24.03$ & $2.73$ & $0.11\times$ & \checkmark \\
\midrule
GLOBAL & $100$ & $1.75$ & $0.0286$ & $0.02\times$ & \checkmark \\
GLOBAL & $1,000$ & $9.62$ & $0.2528$ & $0.03\times$ & \checkmark \\
GLOBAL & $10,000$ & $111.8$ & $2.81$ & $0.03\times$ & \checkmark \\
\midrule
CONTEXT & $100$ & $6.64$ & $0.0933$ & $0.01\times$ & \checkmark \\
CONTEXT & $1,000$ & $31.41$ & $0.6028$ & $0.02\times$ & \checkmark \\
CONTEXT & $10,000$ & $349.7$ & $6.89$ & $0.02\times$ & \checkmark \\
\midrule
TEMPORAL (bounded) & $100$ & $7.41$ & $0.4467$ & $0.06\times$ & \checkmark \\
TEMPORAL (bounded) & $1,000$ & $40.90$ & $4.75$ & $0.12\times$ & \checkmark \\
TEMPORAL (bounded) & $10,000$ & $393.8$ & $57.27$ & $0.15\times$ & \checkmark \\
\midrule
TEMPORAL (unbounded) & $100$ & $5.64$ & $0.4668$ & $0.08\times$ & \checkmark \\
TEMPORAL (unbounded) & $1,000$ & $33.49$ & $4.87$ & $0.15\times$ & \checkmark \\
TEMPORAL (unbounded) & $10,000$ & $322.3$ & $58.91$ & $0.18\times$ & \checkmark \\
\bottomrule
\end{tabular}
\end{table}


\subsection{External Reference Point: Apache Jena TDB2}
\label{sec:eval:jena}

rdflib is a genuine external implementation, but it is an in-memory,
pure-Python library --- not what a systems reviewer means by a
production RDF store. We additionally ran all five query classes
against a real Apache Jena TDB2 store (v6.2.0), a disk-backed,
independently engineered triple store with its own B+tree indexing and
query optimiser, accessed exactly as an external client would: over
HTTP, via the standard SPARQL~1.1 protocol, through a Fuseki server. The
identical fact stream, identical reification predicates, and identical
SPARQL query text used for the rdflib comparison were reused so the two
external comparisons are directly comparable to each other and to SFDB
(Table~\ref{tab:jena}). Every row returned an identical result count to
SFDB.

% Auto-generated by scripts/generate_tables.py from results/jena_reference.json
% Do not edit by hand.
\begin{table}[htbp]
\centering
\caption{External reference point: a real, disk-backed Apache Jena TDB2
store (via Fuseki, queried over HTTP SPARQL) versus SFDB on the same
fact streams (mean over ten runs). The ratio is SFDB\,/\,Jena; below
$1$ means SFDB is faster. The last column
records whether both systems returned the same number of results. Jena
latencies include HTTP round-trip overhead SFDB's in-process calls do
not pay; see the text for how this cuts against SFDB and cannot explain
the crossover on TEMPORAL.}
\label{tab:jena}
\small
\begin{tabular}{llrrrc}
\toprule
Class & Facts & Jena TDB2 (ms) & SFDB (ms) & Ratio & Counts \\
\midrule
LOOKUP & $100$ & $15.39$ & $0.0231$ & $0.00\times$ & \checkmark \\
LOOKUP & $1,000$ & $16.61$ & $0.2187$ & $0.01\times$ & \checkmark \\
LOOKUP & $10,000$ & $31.04$ & $2.77$ & $0.09\times$ & \checkmark \\
LOOKUP & $100,000$ & $71.59$ & $29.39$ & $0.41\times$ & \checkmark \\
\midrule
GLOBAL & $100$ & $12.39$ & $0.0287$ & $0.00\times$ & \checkmark \\
GLOBAL & $1,000$ & $14.32$ & $0.2484$ & $0.02\times$ & \checkmark \\
GLOBAL & $10,000$ & $46.07$ & $3.29$ & $0.07\times$ & \checkmark \\
GLOBAL & $100,000$ & $331.2$ & $37.69$ & $0.11\times$ & \checkmark \\
\midrule
CONTEXT & $100$ & $15.18$ & $0.0641$ & $0.00\times$ & \checkmark \\
CONTEXT & $1,000$ & $16.90$ & $0.5956$ & $0.04\times$ & \checkmark \\
CONTEXT & $10,000$ & $32.77$ & $7.08$ & $0.22\times$ & \checkmark \\
CONTEXT & $100,000$ & $154.0$ & $76.91$ & $0.50\times$ & \checkmark \\
\midrule
TEMPORAL (bounded) & $100$ & $15.21$ & $0.4312$ & $0.03\times$ & \checkmark \\
TEMPORAL (bounded) & $1,000$ & $14.48$ & $4.86$ & $0.34\times$ & \checkmark \\
TEMPORAL (bounded) & $10,000$ & $20.07$ & $59.08$ & $2.94\times$ & \checkmark \\
TEMPORAL (bounded) & $100,000$ & $113.2$ & $774.9$ & $6.85\times$ & \checkmark \\
\midrule
TEMPORAL (unbounded) & $100$ & $13.87$ & $0.4668$ & $0.03\times$ & \checkmark \\
TEMPORAL (unbounded) & $1,000$ & $14.95$ & $5.12$ & $0.34\times$ & \checkmark \\
TEMPORAL (unbounded) & $10,000$ & $28.77$ & $62.83$ & $2.18\times$ & \checkmark \\
TEMPORAL (unbounded) & $100,000$ & $218.8$ & $842.0$ & $3.85\times$ & \checkmark \\
\bottomrule
\end{tabular}
\end{table}


The result is more informative than either a clean win or a clean loss.
On LOOKUP, GLOBAL, and CONTEXT, SFDB remains faster than Jena at every
scale
measured: the ratio (jena\,/\,SFDB mean latency; above $1$ means SFDB is
faster) ranges \rangelookupjena{} on LOOKUP, \rangeglobaljena{} on
GLOBAL, and \rangecontextjena{} on CONTEXT across all scales, settling at \lookupjenalarge{},
\globaljenalarge{}, and \contextjenalarge{} at the largest scale, \jenalargescale{} facts ---
notable given that Jena's reported latency includes an HTTP round trip
SFDB's in-process call does not pay, which if anything understates
SFDB's advantage on these three classes rather than inflating it. On
TEMPORAL, the picture reverses as scale grows: SFDB is
faster at $100$ and $1\,000$ facts, but Jena is faster from $10\,000$
facts onward, with the same ratio falling to
\temporaljenalarge{} (bounded TEMPORAL) and
\temporalunbjenalarge{} (unbounded TEMPORAL) at \jenalargescale{} facts
--- a ratio below $1$ here means Jena is faster, by a factor of roughly
seven at this scale --- despite paying the same HTTP overhead
that works against it on the other three classes. This is not the storage-layer effect
Section~\ref{sec:eval:control} identified for TEMPORAL
against KG-mem: Jena is disk-backed, HTTP-mediated, and not a Python
dict, so its query engine is winning on its own merits --- most likely a
mature cost-based optimiser and compiled index structures overtaking
SFDB's Python-level candidate filtering once the candidate sets reach
tens of thousands of rows. We report this as a genuine, scale-dependent crossover
rather than smoothing it into either "SFDB wins" or "SFDB loses": the
architecture's advantage is real, robust, and now holds on three of the
four query classes measured, and a mature production engine is a real
competitor on the fourth once the working set is large enough for
its optimiser to matter.

\emph{A retraction.} An earlier revision of this section reported an
eighth methodological finding here: that SFDB's own CONTEXT latency at
$100\,000$ facts showed $\sim\!45\%$ run-to-run variance ($712$\,ms
versus $1\,044$\,ms versus $989$\,ms across three independent process
launches) attributed to unspecified process-level state. That finding
is retracted, not merely superseded: those absolute numbers no longer
reproduce at all, because they were measurements of the bug fixed and
described in the next finding (finding ten, Section~\ref{sec:eval:lubm}) ---
CONTEXT queries at that scale now take
\contextjenalarge{}-scale-relative fractions of a millisecond, not
hundreds of them. We do not know whether the specific $\sim\!45\%$
swing was itself a symptom of the same bug or a genuine, separate
process-level effect riding on top of an already-inflated baseline;
distinguishing those would require re-instrumenting a code path that no
longer exists in its old form, which we have not done. We record the
retraction here, in the section that originally reported the now-wrong
finding, rather than silently drop it, on the same principle the rest
of this paper's self-correction record is built on.

\subsection{External Reference Point: OpenLink Virtuoso}
\label{sec:eval:virtuoso}

A single external comparison leaves open whether a reported crossover is
a property of mature production RDF engines generally or an artefact of
one vendor's specific optimiser. We therefore repeated the identical
methodology against a second, independently engineered production
store: OpenLink Virtuoso (open-source edition, v7.20.3243), a
disk-backed C store with its own cost-based SQL/SPARQL optimiser,
unrelated to Jena's Java B+tree implementation and written by a
different organisation. As with Jena, Virtuoso is accessed exactly as an
external client would, over the standard SPARQL~1.1 protocol: HTTP
Digest-authenticated SPARQL Update for inserts and plain SPARQL Query
for reads. Virtuoso requires every inserted triple to name an explicit
graph (Jena does not), so facts are inserted into one named graph;
reads use the identical, unmodified query text shared with the rdflib
and Jena comparisons by passing Virtuoso's \texttt{default-graph-uri}
HTTP parameter rather than rewriting the query. All five query classes
were run at all four paper scales, with every row returning an
identical result count to SFDB (Table~\ref{tab:virtuoso}).

% Auto-generated by scripts/generate_tables.py from results/virtuoso_reference.json
% Do not edit by hand.
\begin{table}[htbp]
\centering
\caption{External reference point: a real, disk-backed OpenLink
Virtuoso store (open-source edition, queried over HTTP SPARQL) versus
SFDB on the same fact streams (mean over ten runs). The ratio is
SFDB\,/\,Virtuoso; below $1$ means SFDB is faster. The last column
records whether both systems returned the same number of results.
Virtuoso latencies include HTTP round-trip and Digest-authentication
overhead SFDB's in-process calls do not pay; see the text for how this
cuts against SFDB and cannot explain the crossover on
TEMPORAL.}
\label{tab:virtuoso}
\small
\begin{tabular}{llrrrc}
\toprule
Class & Facts & Virtuoso (ms) & SFDB (ms) & Ratio & Counts \\
\midrule
LOOKUP & $100$ & $12.47$ & $0.0233$ & $0.00\times$ & \checkmark \\
LOOKUP & $1,000$ & $15.52$ & $0.2201$ & $0.01\times$ & \checkmark \\
LOOKUP & $10,000$ & $29.88$ & $2.58$ & $0.09\times$ & \checkmark \\
LOOKUP & $100,000$ & $108.3$ & $33.05$ & $0.31\times$ & \checkmark \\
\midrule
GLOBAL & $100$ & $14.00$ & $0.0275$ & $0.00\times$ & \checkmark \\
GLOBAL & $1,000$ & $21.63$ & $0.2434$ & $0.01\times$ & \checkmark \\
GLOBAL & $10,000$ & $65.72$ & $2.71$ & $0.04\times$ & \checkmark \\
GLOBAL & $100,000$ & $535.2$ & $35.50$ & $0.07\times$ & \checkmark \\
\midrule
CONTEXT & $100$ & $12.72$ & $0.0632$ & $0.00\times$ & \checkmark \\
CONTEXT & $1,000$ & $16.35$ & $0.5971$ & $0.04\times$ & \checkmark \\
CONTEXT & $10,000$ & $48.32$ & $6.90$ & $0.14\times$ & \checkmark \\
CONTEXT & $100,000$ & $288.7$ & $78.35$ & $0.27\times$ & \checkmark \\
\midrule
TEMPORAL (bounded) & $100$ & $13.96$ & $0.4369$ & $0.03\times$ & \checkmark \\
TEMPORAL (bounded) & $1,000$ & $14.28$ & $4.75$ & $0.33\times$ & \checkmark \\
TEMPORAL (bounded) & $10,000$ & $24.06$ & $55.85$ & $2.32\times$ & \checkmark \\
TEMPORAL (bounded) & $100,000$ & $105.5$ & $775.0$ & $7.34\times$ & \checkmark \\
\midrule
TEMPORAL (unbounded) & $100$ & $9.31$ & $0.4563$ & $0.05\times$ & \checkmark \\
TEMPORAL (unbounded) & $1,000$ & $10.95$ & $4.97$ & $0.45\times$ & \checkmark \\
TEMPORAL (unbounded) & $10,000$ & $26.54$ & $58.20$ & $2.19\times$ & \checkmark \\
TEMPORAL (unbounded) & $100,000$ & $211.8$ & $854.6$ & $4.04\times$ & \checkmark \\
\bottomrule
\end{tabular}
\end{table}


The result reproduces the Jena finding independently, in both direction
and approximate scale threshold. On LOOKUP, GLOBAL, and CONTEXT, SFDB remains
faster than Virtuoso at every scale measured: the ratio
(virtuoso\,/\,SFDB mean latency; above $1$ means SFDB is faster) ranges
\rangelookupvirtuoso{} on LOOKUP, \rangeglobalvirtuoso{} on GLOBAL, and
\rangecontextvirtuoso{} on CONTEXT
across all scales, settling at \lookupvirtuosolarge{},
\globalvirtuosolarge{}, and \contextvirtuosolarge{} at \jenalargescale{} facts --- again despite
Virtuoso's reported latency including an HTTP round trip and
Digest-authentication overhead SFDB's in-process call does not pay. On
TEMPORAL, the same reversal appears at the same scale
threshold as Jena: SFDB is faster at $100$ and $1\,000$ facts, but
Virtuoso is faster from $10\,000$ facts onward, with the ratio falling to
\temporalvirtuosolarge{} (bounded
TEMPORAL) and \temporalunbvirtuosolarge{} (unbounded TEMPORAL) at
\jenalargescale{} facts --- Virtuoso faster by a factor of roughly six
to nine at this scale, close to Jena's margin on the same class.

That two independently engineered production optimisers, from different
organisations, in different implementation languages, reverse the same
comparison on the same query class at the same scale threshold is
stronger evidence than either comparison alone that the crossover is a
property of what mature cost-based query optimisers do once candidate
sets reach tens of thousands of rows, not a Jena-specific quirk. It does
not settle the question --- two data points are not a general law --- but it moves the
open question posed after the Jena result (Section~\ref{sec:future:jena})
from "does this generalise beyond one vendor" to "does this generalise
beyond two"; Section~\ref{sec:eval:graphdb}, next, pushes it to three.

Building this comparison surfaced a ninth methodological finding,
recorded in full in Section~\ref{sec:discussion:threats}: the SPARQL
query text shared across all three external comparisons contains a
FILTER that compares two RDF literals with \texttt{<} or \texttt{>}
(the TEMPORAL classes' range bounds). Against untyped literals, Virtuoso
silently returned every row unfiltered rather than applying a
spec-consistent comparison, while the identical query text against Jena
and rdflib correctly filtered and agreed with SFDB. We fixed this once,
in the shared query text used by all three comparisons, by wrapping the
compared variable in \texttt{STR(...)} to force an unambiguous string
comparison; all three comparisons --- including Jena and rdflib, which
were not broken --- were re-verified to still agree with SFDB after the
change. Unlike the other eight findings, this one is not a bug in code
this project controls: it is evidence that reusing identical query text
across engines, while a necessary discipline for a fair comparison, is
not sufficient on its own when the engines being compared disagree about
SPARQL's own semantics at the margins.

\subsection{External Reference Point: Ontotext GraphDB}
\label{sec:eval:graphdb}

Two independent confirmations narrow the ``is this Jena-specific''
question considerably, but two data points do not establish a general
law. We ran the identical methodology a third time against Ontotext
GraphDB (free edition, v10.7.3), a disk-backed Java store with its own
RDF4J-based storage layer and query optimiser, independently engineered
from both Jena and Virtuoso. GraphDB needs neither authentication nor an
explicit named graph for inserts --- closer in that respect to Jena than
to Virtuoso --- and its plain (untyped) literal relational comparisons
already behave correctly under the SPARQL~1.1 simple-literal ordering
extension, unlike Virtuoso (Section~\ref{sec:eval:virtuoso}), so no
engine-specific workaround was needed beyond the shared
\texttt{STR(...)} coercion already applied uniformly across all three
comparisons. All five query classes were run at all four paper scales,
with every row returning an identical result count to SFDB
(Table~\ref{tab:graphdb}).

% Auto-generated by scripts/generate_tables.py from results/graphdb_reference.json
% Do not edit by hand.
\begin{table}[htbp]
\centering
\caption{External reference point: a real, disk-backed Ontotext
GraphDB store (free edition, queried over HTTP SPARQL) versus SFDB on
the same fact streams (mean over ten runs). The ratio is
SFDB\,/\,GraphDB; below $1$ means SFDB is faster. The last column
records whether both systems returned the same number of results.
GraphDB latencies include an HTTP round trip SFDB's in-process calls do
not pay; see the text for how this cuts against SFDB and cannot explain
the crossover on TEMPORAL.}
\label{tab:graphdb}
\small
\begin{tabular}{llrrrc}
\toprule
Class & Facts & GraphDB (ms) & SFDB (ms) & Ratio & Counts \\
\midrule
LOOKUP & $100$ & $10.72$ & $0.0231$ & $0.00\times$ & \checkmark \\
LOOKUP & $1,000$ & $14.54$ & $0.2259$ & $0.02\times$ & \checkmark \\
LOOKUP & $10,000$ & $27.83$ & $2.57$ & $0.09\times$ & \checkmark \\
LOOKUP & $100,000$ & $55.80$ & $30.94$ & $0.55\times$ & \checkmark \\
\midrule
GLOBAL & $100$ & $13.92$ & $0.0279$ & $0.00\times$ & \checkmark \\
GLOBAL & $1,000$ & $17.19$ & $0.2436$ & $0.01\times$ & \checkmark \\
GLOBAL & $10,000$ & $40.89$ & $2.67$ & $0.07\times$ & \checkmark \\
GLOBAL & $100,000$ & $236.4$ & $37.55$ & $0.16\times$ & \checkmark \\
\midrule
CONTEXT & $100$ & $15.46$ & $0.0634$ & $0.00\times$ & \checkmark \\
CONTEXT & $1,000$ & $13.90$ & $0.5913$ & $0.04\times$ & \checkmark \\
CONTEXT & $10,000$ & $32.38$ & $9.21$ & $0.28\times$ & \checkmark \\
CONTEXT & $100,000$ & $125.9$ & $77.05$ & $0.61\times$ & \checkmark \\
\midrule
TEMPORAL (bounded) & $100$ & $15.55$ & $0.4548$ & $0.03\times$ & \checkmark \\
TEMPORAL (bounded) & $1,000$ & $17.65$ & $4.77$ & $0.27\times$ & \checkmark \\
TEMPORAL (bounded) & $10,000$ & $31.52$ & $58.07$ & $1.84\times$ & \checkmark \\
TEMPORAL (bounded) & $100,000$ & $128.7$ & $790.2$ & $6.14\times$ & \checkmark \\
\midrule
TEMPORAL (unbounded) & $100$ & $14.11$ & $0.4784$ & $0.03\times$ & \checkmark \\
TEMPORAL (unbounded) & $1,000$ & $24.79$ & $4.91$ & $0.20\times$ & \checkmark \\
TEMPORAL (unbounded) & $10,000$ & $36.49$ & $59.17$ & $1.62\times$ & \checkmark \\
TEMPORAL (unbounded) & $100,000$ & $186.3$ & $837.9$ & $4.50\times$ & \checkmark \\
\bottomrule
\end{tabular}
\end{table}


The result is the same pattern a third time, at the same scale
threshold. On LOOKUP, GLOBAL, and CONTEXT, SFDB remains faster than GraphDB at
every scale: the ratio (graphdb\,/\,SFDB mean latency; above $1$ means
SFDB is faster) ranges \rangelookupgraphdb{} on LOOKUP,
\rangeglobalgraphdb{} on GLOBAL, and \rangecontextgraphdb{} on CONTEXT across all scales, settling at
\lookupgraphdblarge{}, \globalgraphdblarge{}, and \contextgraphdblarge{} at \graphdblargescale{}
facts. On TEMPORAL, GraphDB overtakes SFDB from $10\,000$
facts onward --- the identical scale threshold as Jena and Virtuoso ---
with the ratio falling to
\temporalgraphdblarge{} (bounded TEMPORAL) and
\temporalunbgraphdblarge{} (unbounded TEMPORAL) at \graphdblargescale{}
facts, a roughly five-to-six-fold margin: smaller than Jena's and
Virtuoso's margins on the same class, but the same direction, same
class, same scale
threshold.

Three independently engineered production optimisers, from three
different organisations, in three different implementation languages
and storage architectures, all reverse the same comparison on the same
query class at the same scale threshold. That is no longer weak
evidence for a Jena-specific quirk; it is the pattern one would expect
if the reversal is a property of what mature cost-based query
optimisers generally do once candidate sets reach tens of thousands of
rows, rather than a property of any one engine. We still stop short of
calling this a general law --- three engines are not an exhaustive
survey of production RDF stores --- but the question this paper's own future-work section posed after
the first comparison (``is this Jena's optimiser or mature engines
generally'') now has a much better-supported answer than it did after
either the second comparison or the first. And it is narrower than the
question originally posed: on the strength of three independent
production engines, three of the paper's four query classes
(LOOKUP, GLOBAL, CONTEXT) show no crossover of any kind at any scale
measured, in-house baseline or production store; only TEMPORAL does.

\subsection{Real-World Case Study: Wikidata Office-Holders}
\label{sec:eval:wikidata}

Every result so far uses this project's own synthetic generator
(Section~\ref{sec:eval:setup}), parameterised by the same authors who
built the engine it benchmarks. We additionally ran the identical five
query classes and insert benchmark against \wikidatanumfacts{} genuine
facts sourced from Wikidata's public SPARQL endpoint, not constructed to
fit the design: every national head-of-state or head-of-government
``position held'' statement (property P39) for $198$ distinct offices
across $168$ countries that Wikidata records, together with the P580/P582
start/end date qualifiers where present ($92.6\%$ and $89.0\%$ of facts,
respectively). A single such fact --- who held office, in which
position, in which country, from when to when --- is a genuine instance
of the reification problem this paper opens with: it does not reduce to
a subject-predicate-object triple without being split apart. The
snapshot is cached (\texttt{data/wikidata\_case\_study\_raw.json}) and the
benchmark never queries Wikidata live, so the result is reproducible
independent of Wikidata's data changing after the snapshot was taken;
\texttt{scripts/fetch\_wikidata\_case\_study.py}'s module docstring
records exactly how the position sample was constructed, including two
approaches that did not work (an unfiltered label regex hit a 504
timeout against Wikidata's public endpoint; an unrandomised join was
found to be dominated by a handful of anomalous position entities,
returning as few as ten distinct positions in a thousand-row sample)
before settling on a class-filtered label search followed by a
country-name validation pass. This is the same three-in-house-engine
comparison as the rest of this section (KnowledgeGraph, KnowledgeGraphMem,
SheafDatabase), not a fourth external-engine comparison: Jena and
Virtuoso were not re-run against this dataset.

LOOKUP is anchored on Antonio L\'opez de Santa Anna (Q189145), the
single person appearing in the most facts in the snapshot: eleven
non-consecutive terms as President of Mexico between 1833 and 1855, a
genuinely fragmented, repeated-office-holding pattern real history
produces and a synthetic generator would not naturally construct.
CONTEXT is anchored on Greece, the country with the most recorded facts
($195$). TEMPORAL uses a bounded 1990--2010 window and an unbounded
window from 2000 onward. Every query class returned an identical result
count across all three engines (Table~\ref{tab:wikidata-case-study}).

% Auto-generated by scripts/generate_tables.py from results/wikidata_case_study.json
% Do not edit by hand.
\begin{table}[htbp]
\centering
\caption{Real-world case study: 5,337 genuine Wikidata
``position held'' facts (national heads of state and government, with
recorded start/end dates and country), not this project's synthetic
generator. KG, KG-mem, and SFDB columns are latency in ms (mean over
ten runs $\pm$ 95\% CI
half-width). Ratios are SFDB\,/\,KG and SFDB\,/\,KG-mem; below $1$
means SFDB is faster. The last column records three-way cross-engine
result-count agreement (insert has no query result to compare).}
\label{tab:wikidata-case-study}
\footnotesize
\resizebox{\textwidth}{!}{%
\begin{tabular}{lrrrrrc}
\toprule
Class & KG & KG-mem & SFDB & vs KG & vs KG-mem & Counts \\
\midrule
Insert & $420.6 \pm 17.4$ & $101.6 \pm 6.8$ & $328.0 \pm 9.0$ & $0.78\times$ & $3.23\times$ & --- \\
\midrule
LOOKUP & $0.3728 \pm 0.03$ & $0.1394 \pm 0.010$ & $0.0265 \pm 0.001$ & $0.07\times$ & $0.19\times$ & \checkmark \\
GLOBAL & $138.1 \pm 4.8$ & $55.58 \pm 1.8$ & $1.74 \pm 0.21$ & $0.01\times$ & $0.03\times$ & \checkmark \\
CONTEXT & $23.00 \pm 1.3$ & $14.24 \pm 0.43$ & $0.3966 \pm 0.02$ & $0.02\times$ & $0.03\times$ & \checkmark \\
TEMPORAL (bounded) & $58.45 \pm 1.8$ & $25.51 \pm 1.2$ & $21.73 \pm 0.72$ & $0.37\times$ & $0.85\times$ & \checkmark \\
TEMPORAL (unbounded) & $68.86 \pm 3.1$ & $28.99 \pm 1.1$ & $22.78 \pm 0.27$ & $0.33\times$ & $0.79\times$ & \checkmark \\
\bottomrule
\end{tabular}%
}
\end{table}


LOOKUP, GLOBAL, and CONTEXT reproduce the synthetic-data finding cleanly: SFDB is
\wikidatalookupkg{} the KG baseline's latency and \wikidatalookupkgmem{}
KG-mem's on LOOKUP, \wikidataglobalkg{} and \wikidataglobalkgmem{}
on GLOBAL, and \wikidatacontextkg{} and \wikidatacontextkgmem{} on
CONTEXT (ratio is engine\,/\,SFDB; below $1$ means SFDB
is faster) --- SFDB remains faster than both baselines on all three of
the query shapes that show a robust advantage everywhere else in this
paper, on data it did not choose. Greece's $195$ facts are unevenly
distributed across a genuinely deep, fragmented 19th- and 20th-century
political history very unlike the synthetic generator's balanced context
tree, and the CONTEXT margin survives that shift comfortably --- a
useful confirmation, since an earlier revision of this evaluation
measured Greece's CONTEXT query before finding ten
(Section~\ref{sec:eval:lubm}) and reported the opposite result, SFDB
\emph{losing} to KG-mem here; that reading no longer reproduces, for the
same reason finding ten's other reported crossovers do not, and is
retracted along with them.

TEMPORAL alone does not reproduce the synthetic-data finding, and we
report it as measured rather than smooth it into the cleaner three-class
story above. SFDB beats \emph{both} baselines on this real dataset
(\wikidatatemporalkg{} vs.\ KG, \wikidatatemporalkgmem{} vs.\ KG-mem,
bounded; \wikidatatemporalunbkg{} and \wikidatatemporalunbkgmem{}
unbounded) --- on the synthetic benchmark, the temporal advantage does
not survive the KG-mem control at all
(Section~\ref{sec:discussion:kg-wins}), and against the three external
production stores it reverses outright from $10^4$ facts onward
(Sections~\ref{sec:eval:jena}--\ref{sec:eval:graphdb}). We do not have a
confirmed mechanism for why TEMPORAL performs comparatively well here;
a plausible candidate is that Wikidata's
real date distribution (concentrated in the 19th--21st centuries with a
long earlier tail, rather than the synthetic generator's uniform span)
produces temporal-bucket occupancy the index structures handle
differently, but confirming that would require instrumenting the
candidate-set sizes directly, which we have not done.

We read this as the case study doing its job rather than failing to
confirm the headline result: LOOKUP, GLOBAL, and CONTEXT's advantage is robust to
a real, organically-shaped dataset the authors did not construct ---
including surviving a topology shape (many facts concentrated under one
country, embedded in a corpus otherwise unrelated to it) that turned out
to be exactly the shape that exposed finding ten's bug when tried
elsewhere at greater scale, in Section~\ref{sec:eval:lubm} --- while
TEMPORAL remains genuinely data-shape-sensitive, in SFDB's favour here
and against it elsewhere. A single case study on one domain at one scale
is not proof that this generalises further; it is evidence that three
of the four query classes stay robust and the fourth stays genuinely
mixed, under exactly the kind of distributional shift
Section~\ref{sec:discussion:threats}'s dataset-bias threat named as
untested.

\subsection{Standard Benchmark Workload: LUBM}
\label{sec:eval:lubm}

The synthetic generator (Section~\ref{sec:eval:setup}) builds a
balanced context tree: fixed depth, fixed branching factor, the same
shape at every scale. The Wikidata case study
(Section~\ref{sec:eval:wikidata}) is organic but still, in its context
structure, dominated by one populous country among a long tail of
smaller ones. Neither tests a topology with \emph{many} large, mutually
unrelated top-level siblings. LUBM (the Lehigh University Benchmark,
\texttt{sfdb.datasets.lubm})
does: at its largest scale here, $265$ independently generated
universities, each its own top-level context with $2$--$5$ departments
beneath it. We ran the identical five-class benchmark against LUBM data
scaled by university count ($1$, $3$, $26$, $265$ universities, chosen
to land close to this paper's usual four fact-count scales without
forcing an exact match) --- the original loader carried no temporal
data at all, so it has been extended with a one-semester envelope on
\texttt{takesCourse} and \texttt{teacherOf} facts, additive metadata
that does not change what LUBM's own standard Q1--Q14 query set would
return. Every class returned an identical result count across
KnowledgeGraph, KnowledgeGraphMem, and SheafDatabase at every scale
(Table~\ref{tab:lubm}).

% Auto-generated by scripts/generate_tables.py from results/lubm_benchmark.json
% Do not edit by hand.
\begin{table}[htbp]
\centering
\caption{Standard benchmark workload: LUBM (Lehigh University
Benchmark) data instead of this project's synthetic generator, scaled
by university count (\texttt{scripts/lubm\_benchmark.py}) to land
close to this paper's usual four fact-count scales. Latency in ms
(mean over ten runs). Ratios are SFDB\,/\,KG and SFDB\,/\,KG-mem;
below $1$ means SFDB is faster. The last column records three-way
cross-engine result-count agreement.}
\label{tab:lubm}
\footnotesize
\resizebox{\textwidth}{!}{%
\begin{tabular}{llrrrrrc}
\toprule
Class & Facts & KG (ms) & KG-mem (ms) & SFDB (ms) & vs KG & vs KG-mem & Counts \\
\midrule
LOOKUP & $160$ & $0.1509$ & $0.0547$ & $0.0139$ & $0.09\times$ & $0.25\times$ & \checkmark \\
LOOKUP & $1,024$ & $0.1596$ & $0.0637$ & $0.0148$ & $0.09\times$ & $0.23\times$ & \checkmark \\
LOOKUP & $10,429$ & $0.2517$ & $0.0960$ & $0.0282$ & $0.11\times$ & $0.29\times$ & \checkmark \\
LOOKUP & $97,289$ & $0.2550$ & $0.1172$ & $0.0432$ & $0.17\times$ & $0.37\times$ & \checkmark \\
\midrule
GLOBAL & $160$ & $2.58$ & $0.8207$ & $0.0613$ & $0.02\times$ & $0.07\times$ & \checkmark \\
GLOBAL & $1,024$ & $19.58$ & $7.30$ & $0.2832$ & $0.01\times$ & $0.04\times$ & \checkmark \\
GLOBAL & $10,429$ & $217.2$ & $90.69$ & $3.22$ & $0.01\times$ & $0.04\times$ & \checkmark \\
GLOBAL & $97,289$ & $2,263.4$ & $1,133.2$ & $49.89$ & $0.02\times$ & $0.04\times$ & \checkmark \\
\midrule
CONTEXT & $160$ & $1.20$ & $0.3607$ & $0.1269$ & $0.11\times$ & $0.35\times$ & \checkmark \\
CONTEXT & $1,024$ & $2.05$ & $0.3946$ & $0.1438$ & $0.07\times$ & $0.36\times$ & \checkmark \\
CONTEXT & $10,429$ & $9.25$ & $0.5379$ & $0.1926$ & $0.02\times$ & $0.36\times$ & \checkmark \\
CONTEXT & $97,289$ & $71.10$ & $0.7782$ & $0.2139$ & $0.00\times$ & $0.27\times$ & \checkmark \\
\midrule
TEMPORAL (bounded) & $160$ & $0.3178$ & $0.1091$ & $0.5230$ & $1.65\times$ & $4.79\times$ & \checkmark \\
TEMPORAL (bounded) & $1,024$ & $1.50$ & $0.3983$ & $3.27$ & $2.18\times$ & $8.21\times$ & \checkmark \\
TEMPORAL (bounded) & $10,429$ & $37.80$ & $24.53$ & $36.66$ & $0.97\times$ & $1.49\times$ & \checkmark \\
TEMPORAL (bounded) & $97,289$ & $565.0$ & $459.8$ & $471.5$ & $0.83\times$ & $1.03\times$ & \checkmark \\
\midrule
TEMPORAL (unbounded) & $160$ & $0.2798$ & $0.0826$ & $0.4773$ & $1.71\times$ & $5.78\times$ & \checkmark \\
TEMPORAL (unbounded) & $1,024$ & $1.68$ & $0.6618$ & $3.28$ & $1.95\times$ & $4.96\times$ & \checkmark \\
TEMPORAL (unbounded) & $10,429$ & $44.81$ & $27.98$ & $36.27$ & $0.81\times$ & $1.30\times$ & \checkmark \\
TEMPORAL (unbounded) & $97,289$ & $608.2$ & $480.4$ & $508.9$ & $0.84\times$ & $1.06\times$ & \checkmark \\
\bottomrule
\end{tabular}%
}
\end{table}


The first run of this comparison surfaced a genuine bug, not a data-shape
finding, and this is the tenth methodological finding recorded in
Section~\ref{sec:discussion:threats}. CONTEXT is anchored on a single,
fixed department ($68$ facts, every scale); its latency should therefore
be flat as unrelated universities are added elsewhere, and
Section~\ref{sec:eval:control}'s complexity claim already says exactly
that --- \S\ref{sec:complexity} describes CONTEXT as a single $O(1)$
open-set lookup followed by an $O(\text{result size})$ pass, independent
of total corpus size. It was not flat. SFDB's CONTEXT latency at the
smallest LUBM scale ($160$ facts) was $0.68$\,ms; at the largest
($97\,289$ facts, the same $68$-fact result), it was $350$\,ms, a
$515\times$ slowdown driven entirely by unrelated data. The root cause
(\texttt{src/sfdb/sheaf/optimizer.py}, \texttt{\_classify\_semilocal}):
CONTEXT query resolution found every fact whose \emph{exact} context
matched the query and then unioned in \emph{every} open set each of
those facts belonged to, including the ancestor open set named for
\texttt{context:world} --- which every fact in the corpus belongs to,
by construction, whenever the query context is not the root. The direct
$O(1)$ lookup the architecture already supports (the open set literally
named \texttt{context:<c>} already \emph{is} the pre-materialised
sub-tree union, per \texttt{engine.\_assign\_open\_sets}) was never
taken. Final results were always correct --- a downstream filter narrows
the over-broad candidate set back down before it reaches the caller,
which is exactly why cross-engine verification never caught it, and why
the balanced synthetic generator never exposed it either: growing that
generator's corpus also grows the query's own target subtree in
lockstep, so the wasted work and the legitimate work scale together and
look like the same line. LUBM's many-large-unrelated-siblings shape is
what decouples them. The fix makes CONTEXT resolution target the
\texttt{context:<c>} open set directly; \S\ref{sec:complexity}'s
complexity claim did not need to change, the implementation needed to
catch up to it. \texttt{tests/sheaf/test\_context\_query\_scope.py} pins
the fix at the optimiser level (the wasted work was invisible in query
\emph{results}, only in how much internal work produced them, so the
regression test checks the classifier's target-open-set list directly,
confirmed to fail against the pre-fix code and pass against the
post-fix code).

This is not a narrow fix. CONTEXT queries appear in every comparison
this paper reports, and the bug had been present, at whatever severity
each topology's shape happened to trigger, for the entire evaluation.
Every CONTEXT number in this paper --- Table~\ref{tab:query},
Table~\ref{tab:jena}, Table~\ref{tab:virtuoso}, Table~\ref{tab:graphdb},
Table~\ref{tab:wikidata-case-study} --- was re-measured after the fix,
not left as it stood; Section~\ref{sec:eval:jena}'s retraction of its
own earlier ``methodological finding'' about CONTEXT measurement
variance, and Section~\ref{sec:eval:wikidata}'s retraction of a
CONTEXT-loses-to-KG-mem reading, are direct consequences. The corrected
picture is, if anything, a stronger result for the design: SFDB now
shows a robust advantage on three of five query classes
(LOOKUP, GLOBAL, CONTEXT) rather than two, against every baseline and
every external store this paper measures, and the crossover this paper
built its central honesty narrative around narrows to TEMPORAL alone.
We report the bug, the fix, and the corrected numbers together rather
than quietly substituting the corrected numbers and leaving the
narrative as though CONTEXT had always been this fast, because the
finding that the evaluation's own topology diversity is what caught a
real implementation gap is itself evidence for why the diversity effort
(three synthetic-adjacent comparisons, three production stores, one
organic dataset, one standard benchmark) was worth making.

On the corrected numbers: LOOKUP, GLOBAL, and CONTEXT all show SFDB
faster than both baselines at every scale, exactly as in every other
comparison in this paper (\rangelookuplubmkg{} and
\rangelookuplubmkgmem{} on LOOKUP; \rangeloballubmkg{} and
\rangeglobalubmkgmem{} on GLOBAL; \rangecontextlubmkg{} and
\rangecontextlubmkgmem{} on CONTEXT, settling at \contextlubmkglarge{}
and \contextlubmkgmemlarge{} at \lubmlargescale{} facts). TEMPORAL is
the one class that does not show a clean win here either: SFDB is
roughly at parity with both baselines at the two largest scales
(\temporallubmkglarge{} vs.\ KG, \temporallubmkgmemlarge{} vs.\ KG-mem
at \lubmlargescale{} facts, bounded; \temporalunblubmkglarge{} and
\temporalunblubmkgmemlarge{} unbounded), neither a clear win nor the
sharp loss seen against KG-mem on the synthetic benchmark
(Section~\ref{sec:eval:control}) or against the three production
stores. A third dataset, a third shape, and TEMPORAL is the one class
that keeps landing somewhere different every time --- which is itself
consistent with the rest of this paper's account of TEMPORAL as the
query shape genuinely most sensitive to a dataset's actual structure,
rather than a settled win or loss in either direction.

\subsubsection{LUBM Against External Stores: Varying Both Axes at Once}
\label{sec:eval:lubm-external}

Every other comparison in this paper varies exactly one of the two
axes that could explain SFDB's advantage: either the dataset changes
and the comparison engine stays this project's own KG/KG-mem baselines
(the LUBM results above, the Wikidata case study), or the comparison
engine changes and the dataset stays the synthetic generator
(Sections~\ref{sec:eval:jena}--\ref{sec:eval:graphdb}). We ran the
identical LUBM data and query set from this section against all three
external stores as well, so that both axes vary at once: an organic,
many-large-unrelated-siblings topology neither store's optimiser has
seen before, measured through the same HTTP/SPARQL path already used
for the synthetic-generator comparisons.

% Auto-generated by scripts/generate_tables.py
% from results/lubm_jena_reference.json. Do not edit by hand.
\begin{table}[htbp]
\centering
\caption{LUBM data against a real, disk-backed Jena TDB2 store ---
varying both the dataset (LUBM rather than this project's synthetic
generator) and the comparison engine (a real production store rather
than the in-house KG/KG-mem baselines) at once. The ratio is
SFDB\,/\,Jena TDB2; below $1$ means SFDB is faster. The last column
records whether both systems returned the same number of results.}
\label{tab:lubm-jena}
\small
\begin{tabular}{llrrrc}
\toprule
Class & Facts & Jena TDB2 (ms) & SFDB (ms) & Ratio & Counts \\
\midrule
LOOKUP & $160$ & $15.11$ & $0.0050$ & $0.00\times$ & \checkmark \\
LOOKUP & $1,024$ & $11.30$ & $0.0050$ & $0.00\times$ & \checkmark \\
LOOKUP & $10,429$ & $8.82$ & $0.0078$ & $0.00\times$ & \checkmark \\
LOOKUP & $97,289$ & $14.29$ & $0.0055$ & $0.00\times$ & \checkmark \\
\midrule
GLOBAL & $160$ & $15.41$ & $0.0419$ & $0.00\times$ & \checkmark \\
GLOBAL & $1,024$ & $16.83$ & $0.2525$ & $0.02\times$ & \checkmark \\
GLOBAL & $10,429$ & $50.12$ & $4.39$ & $0.09\times$ & \checkmark \\
GLOBAL & $97,289$ & $333.1$ & $38.64$ & $0.12\times$ & \checkmark \\
\midrule
CONTEXT & $160$ & $12.10$ & $0.1061$ & $0.01\times$ & \checkmark \\
CONTEXT & $1,024$ & $9.54$ & $0.1052$ & $0.01\times$ & \checkmark \\
CONTEXT & $10,429$ & $7.91$ & $0.1888$ & $0.02\times$ & \checkmark \\
CONTEXT & $97,289$ & $28.44$ & $0.2088$ & $0.01\times$ & \checkmark \\
\midrule
TEMPORAL (bounded) & $160$ & $12.69$ & $0.5802$ & $0.05\times$ & \checkmark \\
TEMPORAL (bounded) & $1,024$ & $8.10$ & $4.07$ & $0.50\times$ & \checkmark \\
TEMPORAL (bounded) & $10,429$ & $18.73$ & $45.97$ & $2.46\times$ & \checkmark \\
TEMPORAL (bounded) & $97,289$ & $50.90$ & $511.5$ & $10\times$ & \checkmark \\
\midrule
TEMPORAL (unbounded) & $160$ & $12.29$ & $0.4657$ & $0.04\times$ & \checkmark \\
TEMPORAL (unbounded) & $1,024$ & $9.38$ & $3.88$ & $0.41\times$ & \checkmark \\
TEMPORAL (unbounded) & $10,429$ & $21.87$ & $46.95$ & $2.15\times$ & \checkmark \\
TEMPORAL (unbounded) & $97,289$ & $71.04$ & $514.2$ & $7.24\times$ & \checkmark \\
\bottomrule
\end{tabular}
\end{table}

% Auto-generated by scripts/generate_tables.py
% from results/lubm_virtuoso_reference.json. Do not edit by hand.
\begin{table}[htbp]
\centering
\caption{LUBM data against a real, disk-backed Virtuoso store ---
varying both the dataset (LUBM rather than this project's synthetic
generator) and the comparison engine (a real production store rather
than the in-house KG/KG-mem baselines) at once. The ratio is
SFDB\,/\,Virtuoso; below $1$ means SFDB is faster. The last column
records whether both systems returned the same number of results.}
\label{tab:lubm-virtuoso}
\small
\begin{tabular}{llrrrc}
\toprule
Class & Facts & Virtuoso (ms) & SFDB (ms) & Ratio & Counts \\
\midrule
LOOKUP & $160$ & $9.68$ & $0.0049$ & $0.00\times$ & \checkmark \\
LOOKUP & $1,024$ & $5.10$ & $0.0075$ & $0.00\times$ & \checkmark \\
LOOKUP & $10,429$ & $7.30$ & $0.0046$ & $0.00\times$ & \checkmark \\
LOOKUP & $97,289$ & $8.48$ & $0.0047$ & $0.00\times$ & \checkmark \\
\midrule
GLOBAL & $160$ & $12.72$ & $0.0430$ & $0.00\times$ & \checkmark \\
GLOBAL & $1,024$ & $23.28$ & $0.3592$ & $0.02\times$ & \checkmark \\
GLOBAL & $10,429$ & $69.33$ & $2.88$ & $0.04\times$ & \checkmark \\
GLOBAL & $97,289$ & $517.0$ & $39.53$ & $0.08\times$ & \checkmark \\
\midrule
CONTEXT & $160$ & $12.33$ & $0.1130$ & $0.01\times$ & \checkmark \\
CONTEXT & $1,024$ & $5.47$ & $0.1566$ & $0.03\times$ & \checkmark \\
CONTEXT & $10,429$ & $23.14$ & $0.1189$ & $0.01\times$ & \checkmark \\
CONTEXT & $97,289$ & $72.84$ & $0.2784$ & $0.00\times$ & \checkmark \\
\midrule
TEMPORAL (bounded) & $160$ & $11.23$ & $0.4994$ & $0.04\times$ & \checkmark \\
TEMPORAL (bounded) & $1,024$ & $12.37$ & $4.19$ & $0.34\times$ & \checkmark \\
TEMPORAL (bounded) & $10,429$ & $8.74$ & $37.96$ & $4.34\times$ & \checkmark \\
TEMPORAL (bounded) & $97,289$ & $44.29$ & $439.9$ & $9.93\times$ & \checkmark \\
\midrule
TEMPORAL (unbounded) & $160$ & $13.59$ & $0.4989$ & $0.04\times$ & \checkmark \\
TEMPORAL (unbounded) & $1,024$ & $7.71$ & $3.98$ & $0.52\times$ & \checkmark \\
TEMPORAL (unbounded) & $10,429$ & $19.33$ & $35.04$ & $1.81\times$ & \checkmark \\
TEMPORAL (unbounded) & $97,289$ & $49.99$ & $455.6$ & $9.11\times$ & \checkmark \\
\bottomrule
\end{tabular}
\end{table}

% Auto-generated by scripts/generate_tables.py
% from results/lubm_graphdb_reference.json. Do not edit by hand.
\begin{table}[htbp]
\centering
\caption{LUBM data against a real, disk-backed GraphDB store ---
varying both the dataset (LUBM rather than this project's synthetic
generator) and the comparison engine (a real production store rather
than the in-house KG/KG-mem baselines) at once. The ratio is
SFDB\,/\,GraphDB; below $1$ means SFDB is faster. The last column
records whether both systems returned the same number of results.}
\label{tab:lubm-graphdb}
\small
\begin{tabular}{llrrrc}
\toprule
Class & Facts & GraphDB (ms) & SFDB (ms) & Ratio & Counts \\
\midrule
LOOKUP & $160$ & $14.33$ & $0.0051$ & $0.00\times$ & \checkmark \\
LOOKUP & $1,024$ & $12.83$ & $0.0050$ & $0.00\times$ & \checkmark \\
LOOKUP & $10,429$ & $7.68$ & $0.0066$ & $0.00\times$ & \checkmark \\
LOOKUP & $97,289$ & $12.46$ & $0.0076$ & $0.00\times$ & \checkmark \\
\midrule
GLOBAL & $160$ & $9.46$ & $0.0470$ & $0.00\times$ & \checkmark \\
GLOBAL & $1,024$ & $11.15$ & $0.2537$ & $0.02\times$ & \checkmark \\
GLOBAL & $10,429$ & $38.87$ & $4.00$ & $0.10\times$ & \checkmark \\
GLOBAL & $97,289$ & $231.5$ & $44.13$ & $0.19\times$ & \checkmark \\
\midrule
CONTEXT & $160$ & $11.21$ & $0.1130$ & $0.01\times$ & \checkmark \\
CONTEXT & $1,024$ & $12.12$ & $0.1070$ & $0.01\times$ & \checkmark \\
CONTEXT & $10,429$ & $13.65$ & $0.1370$ & $0.01\times$ & \checkmark \\
CONTEXT & $97,289$ & $34.21$ & $0.2713$ & $0.01\times$ & \checkmark \\
\midrule
TEMPORAL (bounded) & $160$ & $12.41$ & $0.5084$ & $0.04\times$ & \checkmark \\
TEMPORAL (bounded) & $1,024$ & $15.62$ & $3.13$ & $0.20\times$ & \checkmark \\
TEMPORAL (bounded) & $10,429$ & $16.68$ & $43.58$ & $2.61\times$ & \checkmark \\
TEMPORAL (bounded) & $97,289$ & $53.69$ & $556.8$ & $10\times$ & \checkmark \\
\midrule
TEMPORAL (unbounded) & $160$ & $11.31$ & $0.4778$ & $0.04\times$ & \checkmark \\
TEMPORAL (unbounded) & $1,024$ & $13.75$ & $3.09$ & $0.22\times$ & \checkmark \\
TEMPORAL (unbounded) & $10,429$ & $27.25$ & $43.57$ & $1.60\times$ & \checkmark \\
TEMPORAL (unbounded) & $97,289$ & $53.24$ & $545.4$ & $10\times$ & \checkmark \\
\bottomrule
\end{tabular}
\end{table}


The pattern reported for the synthetic generator holds exactly, against
every store, at every scale, with every result count matching: SFDB
faster on LOOKUP, GLOBAL, and CONTEXT
(\rangelookuplubmjena{}/\rangelookuplubmvirtuoso{}/\rangelookuplubmgraphdb{}
on LOOKUP; \rangegloballubmjena{}/\rangegloballubmvirtuoso{}/\rangegloballubmgraphdb{}
on GLOBAL; \rangecontextlubmjena{}/\rangecontextlubmvirtuoso{}/\rangecontextlubmgraphdb{}
on CONTEXT, settling at \contextlubmjenalarge{}, \contextlubmvirtuosolarge{},
and \contextlubmgraphdblarge{} respectively at \lubmjenalargescale{} facts),
and TEMPORAL crossing over to favour the external store at the largest
scales (\rangetemporallubmjena{}/\rangetemporallubmvirtuoso{}/\rangetemporallubmgraphdb{}
bounded; \rangetemporalunblubmjena{}/\rangetemporalunblubmvirtuoso{}/\rangetemporalunblubmgraphdb{}
unbounded). This is the strongest robustness check this paper runs for
its central crossover claim: the same three-classes-robust,
one-class-crossover shape survives changing the dataset and the engine
simultaneously, against three independently engineered production
stores, rather than resting on either axis varying alone.

\subsection{Memory}
\label{sec:eval:memory}

The latency results above are not free. Table~\ref{tab:memory} reports the
resident memory delta from inserting $N$ facts into an empty engine,
averaged over ten independent runs, for all three engines.

% Auto-generated by scripts/generate_tables.py from results/paper_suite_summary.json
% Do not edit by hand.
\begin{table}[htbp]
\centering
\caption{Resident memory delta from inserting $N$ facts (mean over ten
runs). Ratios are SFDB\,/\,KG and SFDB\,/\,KG-mem. A dagger ($\dagger$)
marks a ratio computed from at least one delta below the $10$\,MB
measurement floor; these are noise, not signal (repeated measurement
showed them landing on either side of parity across independent runs),
and are excluded from
the ranges reported in the text.}
\label{tab:memory}
\begin{tabular}{lrrrrr}
\toprule
Facts & KG (MB) & KG-mem (MB) & SFDB (MB) & vs KG & vs KG-mem \\
\midrule
$100$ & $0.0$ & $0.0$ & $0.0$ & $0.11\times^\dagger$ & $0.35\times^\dagger$ \\
$1,000$ & $2.1$ & $1.0$ & $1.7$ & $0.82\times^\dagger$ & $1.80\times^\dagger$ \\
$10,000$ & $26.7$ & $25.1$ & $47.6$ & $1.78\times$ & $1.90\times$ \\
$100,000$ & $285.2$ & $376.6$ & $703.9$ & $2.47\times$ & $1.87\times$ \\
\bottomrule
\end{tabular}
\end{table}


\begin{figure}[htbp]
\centering
\includegraphics[width=0.6\textwidth]{figures/memory.pdf}
\caption{Resident memory delta from inserting $N$ facts, all three
engines. Ratios are computed only at scales where both deltas exceed
the measurement floor; see text.}
\label{fig:memory}
\end{figure}

SFDB uses \rangememkg{} more resident memory per fact than the KG
baseline, and \rangememkgmem{} more than the KG-mem control, over the
scales at which both deltas are large enough to measure reliably
($\geq 1$\,MB); at $100$ facts the deltas sit at the measurement floor
and are excluded from the ranges rather than reported as signal. The
cost has a direct structural explanation: the topology
builder assigns each fact to roughly eight to ten open sets (event,
subject and object entities, temporal bucket, context and its ancestors,
and two provenance buckets), and the stalk, neighbourhood, context, and
temporal indexes each hold their own reference to the fact. The KG
baseline's SQLite-backed triple table is more compact per fact because
reified triples share dictionary-encoded integer identifiers rather than
holding live Python object references; KG-mem gives some of that
compactness back by holding its indexes in Python dicts, which is why
SFDB's multiple against it is smaller. This is the paper's central
tradeoff: SFDB spends memory on redundant indexing structure to buy the
latency advantage reported above.

\subsection{Scalability}
\label{sec:eval:scalability}

Figure~\ref{fig:scalability} plots latency against dataset size on log axes.
Insert latency scales linearly for both write-through engines. SFDB LOOKUP
and GLOBAL both scale linearly in the number of facts, since both resolve
through the same flat, hash-indexed stalk table; the residual gap to KG is
a constant factor set by the cost of SQLite round trips and
reified-triple reconstruction rather than a difference in asymptotic
complexity, and the gap to KG-mem is the reconstruction share alone.
SFDB's unbounded TEMPORAL queries scale as $O(\log M + n)$ in the number
of temporally-annotated facts $M$ (a binary search into the flat
start/end index plus the candidate set $n$); KG's TEMPORAL cost is
$O(T)$ regardless of whether the range is bounded, since it always scans
and filters every reified temporal triple.

\begin{figure}[htbp]
\centering
\includegraphics[width=0.6\textwidth]{figures/scalability.pdf}
\caption{Scalability: latency versus dataset size (log-log), KG baseline
and SFDB.}
\label{fig:scalability}
\end{figure}

\subsection{Cross-Engine Verification}
\label{sec:eval:verification}

Every benchmark row above was accompanied by a semantic equivalence check
across all three engines.
The check canonicalises each engine's result set and compares them as
order-independent sets of canonical facts, so a reordering of otherwise
identical results does not register as a divergence and a genuine content
difference cannot hide behind row order. Over the full evaluation
(\verifiedtotal{} query-class-by-scale combinations, each backed by ten
timed runs) we observed \verifiedcount{} verified and zero unresolved
divergences. This does not, on its own, prove
Theorem~\ref{thm:equivalence}; the theorem is proved constructively in
Section~\ref{sec:correctness}. It does, however, provide direct empirical
corroboration on every query and input reported in this section, and it is
the same harness that caught the LOOKUP semantics bug described in
Section~\ref{sec:correctness}. The rdflib comparison adds a
result-count check against an implementation we did not write.

\subsection{Summary}
\label{sec:eval:summary}

Six findings summarise the evaluation:

\begin{enumerate}
    \item \textbf{Inserts do not favour SFDB unconditionally}: SFDB is
          $1.2\times$ \emph{slower} than the write-through KG baseline
          at the smallest scale ($100$ facts) and $1.3$--$1.5\times$
          faster at every larger scale measured
          (\rangeinsertkg{} across the full range). The reversal at
          small $N$ is consistent with the fixed cost of registering a
          fact into eight to ten open sets not yet being amortised.
          (KG-mem inserts faster than both by
          skipping persistence; that row is a control artefact, not a
          finding.)

    \item \textbf{LOOKUP queries are \rangelookupkg{} faster on SFDB
          than the KG baseline, and \rangelookupkgmem{} faster than the
          KG-mem control}: the bulk of the advantage survives removal of
          the storage-layer difference and is attributable to returning
          stored facts instead of reassembling them from reified
          triples.

    \item \textbf{GLOBAL queries are \rangeglobalkg{} faster than KG
          and \rangeglobalkgmem{} faster than KG-mem}, for the same
          reason: the flat stalk index yields whole facts; both KG
          variants pay per-fact reconstruction.

    \item \textbf{CONTEXT queries --- the class the context-indexed
          design is named for, and the one no earlier revision of this
          evaluation measured --- are \rangecontextkg{} faster than KG
          and \rangecontextkgmem{} faster than KG-mem}: the anchor
          context's sub-tree is pre-materialised as a single open set at
          insert time, so the query resolves without a scan over the
          context space.

    \item \textbf{The TEMPORAL advantages over the SQLite baseline
          (\rangetemporalkg{} bounded, \rangetemporalunbkg{} unbounded)
          do not survive the control --- and on bounded TEMPORAL, SFDB is
          slower than KG-mem at three of four scales}
          (\rangetemporalkgmem{} bounded, \rangetemporalunbkgmem{}
          unbounded, both straddling parity): these are storage-layer
          effects, not wins for context-indexed
          storage, and once controlled for there is no consistent
          advantage left to report on either TEMPORAL class.

    \item \textbf{SFDB uses \rangememkg{} more resident memory per fact
          than KG, and \rangememkgmem{} more than KG-mem}: the latency
          advantages above are funded by keeping redundant per-open-set
          index structure in memory.
\end{enumerate}

We do not claim SFDB is universally superior; the memory cost is real,
inserts are not uniformly faster, the temporal classes show no
design-level advantage once the storage
layer is controlled for and are mildly \emph{disadvantageous} on bounded
TEMPORAL, and all of this is reported plainly rather than smoothed into a
single directional headline.

\subsection{Limitations of the Evaluation}
\label{sec:eval:limitations}

The baseline and scale caveats that qualify every number above are
collected in one place in Section~\ref{sec:limitations} and analysed in
Section~\ref{sec:discussion:threats} rather than repeated here; in short,
the KG baselines are our own research prototypes rather than production
RDF stores, rdflib is the only externally-authored system measured, and
all measurements stop at $10^5$ facts on a single node.
